% ─────────────────────────────────────────────────────────────────
% donkified_math.tex
% The math behind the oracle detector, explained to a donkey.
% Compile: pdflatex donkified_math.tex
% ─────────────────────────────────────────────────────────────────
\documentclass[11pt]{article}

\usepackage[margin=1in]{geometry}
\usepackage{amsmath, amssymb}
\usepackage{xcolor}
\usepackage{mdframed}
\usepackage{enumitem}
\usepackage{hyperref}
\hypersetup{colorlinks=true, linkcolor=brown, urlcolor=brown}

\newcommand{\donkey}[1]{%
  \begin{mdframed}[backgroundcolor=brown!8, linecolor=brown!50, linewidth=1pt,
                   roundcorner=6pt, skipabove=8pt, skipbelow=8pt]
  \textbf{\textcolor{brown}{\large Donkey says:}}\quad #1
  \end{mdframed}}

\title{\textbf{Donkified Math}\\[4pt]
       \large The Oracle Detector, Explained to a Donkey}
\author{HIL Controller Benchmark --- \texttt{oracle\_detector\_node.py}}
\date{\today}

\begin{document}
\maketitle

\begin{abstract}
\noindent
There is a donkey. The donkey is a drone. Somewhere in town there is a
\textbf{carrot} (the target). A normal detector (YOLO) would squint at a
camera photo and \emph{guess} where the carrot is. Our detector --- the
\textbf{oracle} --- does not squint at anything. It \emph{reads the answer key}:
it already knows, as plain numbers, where the donkey stands and where the carrot
sits. Then it does a little bit of arithmetic to work out exactly where the
carrot \emph{would} appear in the donkey's eye, and draws a perfect box around
it. This document is that little bit of arithmetic, donkey-sized.
\end{abstract}

% ─────────────────────────────────────────────────────────────────
\section{The Big Idea (read this even if you read nothing else)}

\donkey{I am not \emph{looking} for the carrot. The town\,(simulator) whispers
the carrot's exact address in my ear. My only job is: ``if the carrot is
\emph{there}, and my head is pointed \emph{this way}, where does it land in my
eyeball, and how big does it look?''}

A real eye (YOLO) goes \textbf{photo $\to$ guess}. The oracle goes
\textbf{known numbers $\to$ exact answer}. No guessing. Same answer every single
time. That is the whole trick, and it is why the oracle is ``cheating'': it
reads coordinates instead of pixels.

The inputs the town hands us for free:
\begin{itemize}[nosep]
  \item \texttt{/sim/drone\_pose} $=(x,y,z,\text{pitch},\text{yaw})$ --- where the donkey is and where its nose points.
  \item \texttt{/sim/target\_pose} $=(t_x,t_y,t_z,\text{yaw})$ --- where the carrot is.
\end{itemize}
The output we must produce: a box on the eyeball, $(u,v,b_w,b_h)$ ---
centre pixel $(u,v)$ and box size width$\times$height.

% ─────────────────────────────────────────────────────────────────
\section{The Cast of Symbols}

\begin{center}
\begin{tabular}{cl}
\hline
Symbol & What it is (in donkey) \\
\hline
$(x,y,z)$ & where the donkey stands, in town coordinates \\
$(t_x,t_y,t_z)$ & where the carrot sits, in town coordinates \\
$\text{yaw}$ & which way the donkey's nose points (spin on the spot) \\
$\text{pitch}$ & how much the donkey nods its head up/down \\
$f_x, f_y$ & ``eyeball zoom'' (focal length), $=554$ \\
$(c_{x0}, c_{y0})$ & the middle of the eyeball, $=(320,240)$ \\
$H$ & the carrot's real height in metres, $=1.5$ (we just \emph{know} it) \\
$Z_c$ & how far ahead the carrot is, straight out the nose \\
$(u,v)$ & where the carrot lands on the eyeball, in pixels \\
$(b_w,b_h)$ & how wide/tall the carrot's box looks \\
\hline
\end{tabular}
\end{center}

% ─────────────────────────────────────────────────────────────────
\section{Step 1 --- ``How far is the carrot from my nose?''}

First, point an arrow from the donkey to the carrot, in plain town coordinates.
That is just subtraction:
\begin{align}
d_x &= t_x - x, & d_y &= t_y - y, & d_z &= t_z - z.
\end{align}

\donkey{Carrot minus me. If the carrot is at town-corner 35 and I am at
corner 10, the carrot is ``25 east of my hooves.'' Easy.}

But this arrow is written in \emph{town} directions (north, east, up). The donkey
does not think in north/east. The donkey thinks in \emph{my-front, my-left,
my-up}. So we must re-write the arrow from the donkey's point of view.

% ─────────────────────────────────────────────────────────────────
\section{Step 2 --- ``Turn the map to face my nose'' (yaw)}

The donkey is spun around by an angle $\text{yaw}$. To turn the town-arrow into a
me-arrow, we rotate it. This is the famous ``spin the map until it lines up with
where I'm walking'' move:
\begin{align}
\text{fwd}  &= d_x\cos(\text{yaw}) + d_y\sin(\text{yaw}), \\
\text{left} &= -d_x\sin(\text{yaw}) + d_y\cos(\text{yaw}), \\
\text{up}   &= d_z .
\end{align}

\donkey{If I am facing backwards (yaw $=\pi$), then a carrot that is ``ahead in
town'' is actually \emph{behind me}. The sines and cosines do exactly this
bookkeeping so I don't fool myself. \texttt{fwd} now means truly in front of
\emph{my} face.}

Now the arrow is in donkey-body talk: how far \emph{forward}, how far to my
\emph{left}, how far \emph{up}.

% ─────────────────────────────────────────────────────────────────
\section{Step 3 --- ``Account for my nodding head'' (pitch)}

The donkey can also nod its head down (to see the ground) or up. That is
\emph{pitch}. Nodding tilts ``forward'' and ``up'' into each other. With
$\theta = (\text{pitch\_sign})\cdot\text{pitch}$:
\begin{align}
\text{fwd}_p &= \text{fwd}\cos\theta + \text{up}\sin\theta, \\
\text{up}_p  &= -\text{fwd}\sin\theta + \text{up}\cos\theta .
\end{align}

\donkey{When I nod down to look at the floor, something that was ``straight
ahead'' is now ``up and ahead'' from my tilted eyeball's view. This step does my
neck for me.}

% ─────────────────────────────────────────────────────────────────
\section{Step 4 --- ``Rename things the way the eyeball likes''}

The donkey's body says \emph{forward / left / up}. But a camera eyeball is
stubborn and insists on its own naming: \emph{$Z_c$ forward, $X_c$ to the
right, $Y_c$ downward}. So we just relabel (note left becomes $-$right, up
becomes $-$down):
\begin{align}
Z_c &= \text{fwd}_p, & X_c &= -\,\text{left}, & Y_c &= -\,\text{up}_p .
\end{align}

\donkey{Same arrow, new sticky-labels. The eyeball is left-handed and grumpy. I
humour it. $Z_c$ is the only one that really matters emotionally: it is ``how far
out my nose the carrot is.''}

% ─────────────────────────────────────────────────────────────────
\section{Step 5 --- ``Far things look small'' (the pinhole)}

This is the heart of seeing. To put a 3D carrot onto a flat eyeball, you
\textbf{divide by how far away it is}. That single division is why distant carrots
look tiny and near carrots look huge:
\begin{align}
u &= f_x\,\frac{X_c}{Z_c} + c_{x0}, \\
v &= f_y\,\frac{Y_c}{Z_c} + c_{y0}.
\end{align}

\donkey{Hold a carrot at arm's length: small. Bring it to my nose: fills my whole
eye. Nothing changed about the carrot --- only $Z_c$ (the distance I divide by)
got smaller. Dividing by distance \emph{is} perspective. The $+c_{x0}$ just slides
the origin to the middle of the eyeball, because pixel $(0,0)$ is the corner, not
the centre.}

% ─────────────────────────────────────────────────────────────────
\section{Step 6 --- ``How big is the box?''}

We \emph{know} the carrot is $H = 1.5$\,m tall in real life. A thing of known
height, seen from distance $Z_c$, casts a box of pixel-height:
\begin{align}
b_h &= \frac{f_y\,H}{Z_c}, & b_w &= b_h \cdot \text{aspect}.
\end{align}

\donkey{Same division-by-distance rule. Carrot far ($Z_c$ big) $\Rightarrow$ box
small. Carrot near ($Z_c$ small) $\Rightarrow$ box big. The box \emph{breathes}:
it grows as I approach, shrinks as I back off --- exactly like a real eye, except
mine is never wrong.}

This breathing box is the donkey's only sense of distance: the controller watches
$b_h$ grow and knows ``I'm getting closer.''

% ─────────────────────────────────────────────────────────────────
\section{Step 7 --- ``When do I see nothing?'' (the gates)}

Even an oracle is allowed to say ``no carrot visible.'' It returns nothing when:
\begin{enumerate}[nosep]
  \item \textbf{Behind me:} $Z_c \le 0.1$. You cannot see what is behind your own
        head. (This is why a donkey starting at yaw$=\pi$ must spin around first.)
  \item \textbf{Too far (optional):} $Z_c > \text{max\_range}$. A ``you must walk
        closer before it counts'' rule, off by default.
  \item \textbf{Off the edge of my eye:} $u,v$ fall outside the
        $640\times480$ eyeball.
  \item \textbf{Stale whisper:} the town hasn't updated my own position in
        $0.3$\,s --- I might be remembering an old run, so I shut up rather than
        lie.
\end{enumerate}

\donkey{Saying ``I see nothing'' is itself information: it makes me \emph{search}.
A wrong ``I see it!'' would make me charge a carrot that isn't there.}

% ─────────────────────────────────────────────────────────────────
\section{The Whole Donkey, in One Breath}

Stack all the steps and the carrot's box is one clean recipe:
\[
\underbrace{(t_x,t_y,t_z)}_{\text{carrot address}}
\;\xrightarrow{\;-\,(x,y,z)\;}\;
\underbrace{(d_x,d_y,d_z)}_{\text{arrow in town}}
\;\xrightarrow{\text{yaw}}\;
\xrightarrow{\text{pitch}}\;
\underbrace{(X_c,Y_c,Z_c)}_{\text{arrow in eyeball}}
\;\xrightarrow{\;\div Z_c\;}\;
\underbrace{(u,v,b_w,b_h)}_{\text{perfect box}}.
\]

Three ideas carry the entire thing:
\begin{enumerate}[nosep]
  \item \textbf{Subtract} to get the arrow (carrot minus me).
  \item \textbf{Rotate} the arrow into my own point of view (yaw, then pitch, then eyeball labels).
  \item \textbf{Divide by distance} to flatten 3D into the eye --- and that same
        division makes the box breathe.
\end{enumerate}

% ─────────────────────────────────────────────────────────────────
\section{The Moral (why this matters)}

\donkey{Because I read the answer key, every donkey-run is identical and perfect.
So if the IBVS donkey and the Proportional donkey behave differently, it is
\emph{their steering brains} that differ --- not their eyes. The oracle removes
the eyes from the argument. That is the entire point: a fair race needs every
runner to see the exact same thing.}

And the punchline from our earlier chat: since the oracle already \emph{knows} the
carrot's position \emph{and} which way it faces (it read both off the answer key),
adding SLAM gives the donkey nothing new. SLAM only matters the day we hide the
answer key and force the donkey to figure out its own heading from the scenery ---
like a real animal with real eyes.

\end{document}